\section{Both sides at every curve point}\label{sec:boundary}

Routing alone would allow the possibility that every off-curve point
has bit zero. We must construct bit-one points, and we must do so
arbitrarily near every curve point. This is the remaining geometric
part of the theorem.

\subsection{A small rational square and two parameter half-arcs}

First let $t_0$ be rational modulo one, and write
\[
 c=\gamma(t_0),\qquad d=\gamma(t_0+1/2),\qquad
 D=2^{-V(1)}.
\]
The inverse separation assumption gives $\norm{c-d}\ge D$.
For a requested radius $2^{-n}$, put
\begin{equation}
 m=\max(n+4,V(1)+4,4),\qquad e=2^{-m},\qquad
 q=G(t_0,m+4).
 \label{eq:localsquare}
\end{equation}
The rational center $q$ satisfies $\norm{q-c}\le e/16$.
Let $Q$ be the boundary of the rational square
$q+[-e,e]^2$. We have
\begin{align}
 \dist(c,Q)&\ge15e/16,\label{eq:insideq}\\
 \dist(d,Q)&\ge D-17e/16\ge239e/16,\label{eq:outsideq}\\
 \norm{z-c}&\le17e/16<2^{-n}\quad(z\in Q).\label{eq:qnear}
\end{align}
The point $c$ is strictly inside this square; $d$ is strictly outside.
Thus
\begin{equation}
                      \parity(Q,c)=1,\qquad \parity(Q,d)=0.
 \label{eq:qvalues}
\end{equation}
These parities at named real points are calculated through rational
approximations using their displayed clearances.

Split the parameter circle into the two closed half-intervals with
endpoints $t_0,t_0+1/2$. Let their images be $A$ and $B$.
Each is a uniformly continuous path from $c$ to $d$,
with direction irrelevant to parity.
Set
\begin{equation}
 r=U(m+4),\qquad K=V(r),\qquad \eta=2^{-K}.
 \label{eq:localscale}
\end{equation}
Subdivide every side of $Q$ into
\begin{equation}
                     T=2^{K-m+7}
 \label{eq:Tsquare}
\end{equation}
equal segments. The exponent is nonnegative by the normalization
of the moduli, and every segment has length exactly $\eta/64$.
There are $4T$ segments, all specified before any color decision.

\subsection{A finite color test with a genuine overlap margin}

\begin{lemma}[The two half-arcs cannot both be close]\label{lem:arcs}
For each short square segment $E$,
\[
       \dist(E,A)<\eta/4\quad\Longrightarrow\quad
       \dist(E,B)\ge\eta/4.
\]
\end{lemma}

\begin{proof}
If a point $a\in A$ and a point $b\in B$ were each within
$\eta/4$ of $E$, then
\[
                  \norm{a-b}<\eta/4+\eta/64+\eta/4<\eta.
\]
For rational parameters the inverse condition forces their circular
distance below $2^{-r}$. Two parameters in opposite half-intervals
that are this close must both be within $2^{-r}$ of the same
endpoint, either $t_0$ or $t_0+1/2$. This is an interval calculation:
after translating $t_0$ to zero, the parameters lie in
$[0,1/2]$ and $[1/2,1]$, and the short gap is either across
$1/2$ or across $0=1$.
The forward modulus puts $a$ within $e/16$ of $c$ or $d$.
It follows that the square is within $e/16+\eta/4$ of that
endpoint, contradicting \eqref{eq:insideq} or \eqref{eq:outsideq}.

All positive closeness witnesses used in the algorithm come from
rational parameter samples. The conclusion for an arbitrary named
parameter follows by approximation. Equivalently, the weak lower
bound on the distance to $B$ means its rational approximations
satisfy the corresponding lower bounds at every precision.
If any such bound failed with positive slack, the preceding
finite rational contradiction would apply. No decision of
real equality or real trichotomy is involved.
\end{proof}

Approximate $\dist(E,A)$ by a rational number $D_E$ with error
at most $\eta/256$. To be explicit, use a mesh of the half-interval
with $2^{U(K+9)+1}$ subdivisions and evaluate at precision $K+9$.
Oscillation and evaluation together cost at most
$2\cdot2^{-(K+9)}=\eta/256$. Minimize the exact rational
point-to-segment distances as in Section~\ref{sec:nets}.
Color the segment by
\begin{equation}
 \begin{array}{lll}
 \text{red}&\text{if }D_E\ge\eta/8,&
       \dist(E,A)\ge31\eta/256>\eta/16,\\
 \text{blue}&\text{if }D_E<\eta/8,&
       \dist(E,B)\ge\eta/4.
 \end{array}
 \label{eq:squarecolors}
\end{equation}
For the blue case,
$\dist(E,A)<33\eta/256<\eta/4$, so
Lemma~\ref{lem:arcs} applies. The estimates leave a specified
rational gap on both branches.

A vertex common to segments of different colors has distance at least
$\eta/16$ from $A$ by its red segment, and at least $\eta/4$
from $B$ by its blue segment. It is therefore off the entire
curve with clearance at least $\eta/16$.

The coloring cannot be monochrome. If all segments were red,
the arc $A$ would be separated from $Q$ by $\eta/16$; if all
were blue, the arc $B$ would be separated from $Q$ by $\eta/4$.
In either case that arc joins $c$ to $d$ while avoiding the
closed walk $Q$, so Proposition~\ref{prop:stable} would give
$\parity(Q,c)=\parity(Q,d)$, contradicting \eqref{eq:qvalues}.
Since the $4T$ colors form a finite cyclic list, there are
at least two color-change vertices.

\subsection{The transition vertices cannot all have one curve bit}

List the transition vertices in cyclic order as
$w_0,\ldots,w_{s-1}$, and let $Q_j$ be the monochromatic
run from $w_j$ to $w_{j+1}$, with subscripts modulo $s$.
These are rational polygonal paths along the square, and their
concatenation is $Q$.
Each $w_j$ has the clearance $\eta/16$, so its curve bit is
a bounded calculation.

Suppose all these bits were equal. The connection theorem supplies
for each $j$ a rational polygonal route $P_j$ from $w_j$ to
$w_{j+1}$, with positive clearance from the entire curve.
Let $R$ be the closed walk obtained by concatenating all $P_j$,
and let $L_j$ be the closed walk following $P_j$ and then
the reverse of $Q_j$.

The walk $R$ avoids $A$, and $A$ joins $c$ to $d$; hence
\[
                         \parity(R,c)=\parity(R,d).
\]
If $Q_j$ is red, $L_j$ avoids $A$; if it is blue, $L_j$ avoids
$B$. In both cases the appropriate half-arc joins $c$ to $d$,
so
\[
                         \parity(L_j,c)=\parity(L_j,d).
\]
The finite chain identity modulo two is
\[
                         \sum_{j<s}L_j=R+Q.
\]
Taking parity at $c$ and at $d$, then xoring the preceding
equalities, forces $\parity(Q,c)=\parity(Q,d)$.
This again contradicts \eqref{eq:qvalues}.

The assertion that all transition bits are equal is decidable.
The contradiction therefore proves that a bounded scan of the
at most $4T$ vertices finds one with bit zero and one with bit one.
Both have clearance at least $\eta/16$ and both are within
$2^{-n}$ of $c$ by \eqref{eq:qnear}.
The auxiliary routes $P_j$ in the contradiction do not need to
stay in the small square or near $c$.

\begin{proposition}[Local witnesses with explicit bounds]\label{prop:local}
For rational $t_0$ and target precision $n$, the construction
\eqref{eq:localsquare}--\eqref{eq:Tsquare} returns rational
points of both curve bits, each within $2^{-n}$ of
$\gamma(t_0)$ and each with clearance at least
$2^{-K-4}=\eta/16$. Its candidate list has at most
$4\cdot2^{K-m+7}$ elements.
\end{proposition}

For a named parameter $t$, first take a rational parameter $t_0$
with $d_\T(t,t_0)\le2^{-U(n+2)}$. Such a parameter is obtained
by reading the parameter name at the specified precision.
Then $\norm{\gamma(t)-\gamma(t_0)}\le2^{-(n+2)}$.
Apply Proposition~\ref{prop:local} with target $n+2$ at
$t_0$. Its two outputs are within $2^{-(n+1)}<2^{-n}$ of
$\gamma(t)$. This proves the both-sides clause for arbitrary
named curve points without choosing an irrational split parameter.

\section{Completing the theorem}\label{sec:completion}

\subsection{The two regions and their sizes}

Define $I$ and $O$ to be the off-curve points with bits one and
zero, respectively. The classification calculation places every
off-curve point in exactly one of them. Local constancy makes
both quantitatively open: a point with clearance $r$ has a
ball of radius $r/4$ of its own bit. The connection theorem
joins any two points of either region by a positive-clearance
polygonal route. Conversely, Proposition~\ref{prop:stable}
prevents such a route between different regions.

If a closed polygonal walk lies in a convex box and the query is
outside that box, its parity is zero: join each edge to the first
vertex to form a triangle fan. Every triangle avoids the query,
and its two radial edges cancel with its neighbors.
Fine approximations of $\Gamma$ lie in the box used in
\eqref{eq:R}; hence every point outside that box has curve bit zero.
If the gap from the box is given by a positive rational number,
the required precision is specified by that gap.
For a merely asserted strict gap, instantiate its positive
rational witness; no unbounded search is required for the proof.

Thus $I\subset[-R,R]^2$. It is inhabited by
Proposition~\ref{prop:local} at $t_0=0$.
The points $(R+j+2,0)$, for $j\in\N$, have bit zero and
clearance at least $j+2$; this proves the asserted constructive
unboundedness of $O$. This also fixes the names: bit one is
inside and bit zero is outside.

\subsection{Approximate intersections of crossing paths}

Let a named uniformly continuous path $\alpha:[0,1]\to\R^2$
have endpoints of different bits, and fix $\eps=2^{-n}$.
Using its modulus and $U$, construct finite rational nets
$X_i$ and $Y_j$ for the path and the curve, respectively,
each with total net error at most $\eps/8$ and individual
evaluation error at most $\eps/8$. These nets include the
appropriate endpoint samples. Compute the rational number
\[
                         D=\min_{i,j}\norm{X_i-Y_j}.
\]
If $D\ge\eps/2$, then for arbitrary path and curve parameters,
the net estimates give distance at least
$D-\eps/8-\eps/8\ge\eps/4$.
The whole path would have positive clearance, contradicting
the endpoint bits and Proposition~\ref{prop:stable}.
The comparison is rational and decidable, so $D<\eps/2$.
A least minimizing pair of indices supplies rational parameters
$s_i,t_j$ with
\[
 \norm{\alpha(s_i)-\gamma(t_j)}
      \le D+\eps/8+\eps/8<3\eps/4<\eps.
\]
The algorithm has a completely specified finite product of meshes
as its search space. It proves approximate contact at each
requested precision. It does not choose an exact zero of a
continuous real function.

\subsection{The exact common-boundary converse by Cauchy names}
\label{sec:boundary-converse}

We have already constructed both bits in every neighborhood of every
curve point. For completeness, we prove the converse in the precise
positive-neighborhood sense stated in Theorem~\ref{thm:main}.
Suppose a named point $x$ has the property that, for every positive
dyadic radius, its neighborhood contains a \emph{rational}
off-curve point of each bit. More precisely, its arithmetic
premise is: for every $n$ there exist rational point codes
$a_0,a_1$ and natural clearance exponents $k_0,k_1$ such that
$\operatorname{Off}(a_i,k_i)$, $b_\Gamma(a_i)=i$, and
$\norm{a_i-x}<2^{-n}$. Strict closeness is expressed by a
finite precision witness in the name of $x$.
These are ordinary number-coded existential witnesses;
a separate choice sequence is not needed.

This rational formulation does not weaken the neighborhood
condition for named points. Given a named off-curve point,
approximate it rationally within a quarter of its positive
clearance and within a quarter of the remaining neighborhood
margin. Local constancy preserves its bit, and the approximation
has a positive clearance. Thus arbitrary named witnesses can
always be replaced by rational ones.

For any dyadic $\tau>0$, sample the curve and $x$ sufficiently
accurately that the combined net and query error is at most
$\tau/8$. Let $D$ be the minimum rational query-to-sample
distance. If $D\ge\tau/2$, every curve point is at least
$3\tau/8$ from $x$. Choose, from the neighborhood premise,
points of both bits within $\tau/16$ of $x$. The segment
joining them lies in that ball and has clearance at least
$3\tau/8-\tau/16>0$. Its endpoint bits must be equal,
a contradiction. Therefore $D<\tau/2$.
The bounded argmin returns a rational parameter $s(\tau)$ such that
\begin{equation}
                     \norm{\gamma(s(\tau))-x}<\tau.
 \label{eq:nearparameter}
\end{equation}
Notice that this parameter is computed just from the point and
curve names. The neighborhood premise is used to prove the
calculation succeeds, not to select an infinite family of data.

Now set
\[
 \tau_n=2^{-V(n+4)-3},\qquad t_n=s(\tau_n).
\]
For $j\ge n$, monotonicity of $V$ and \eqref{eq:nearparameter}
give
\[
 \norm{\gamma(t_j)-\gamma(t_n)}
    <\tau_j+\tau_n\le 2^{-V(n+4)-2}<2^{-V(n+4)}.
\]
The inverse condition, applied to the rational parameters,
therefore yields
\begin{equation}
                         d_\T(t_j,t_n)<2^{-(n+4)}.
 \label{eq:circlecauchy}
\end{equation}
This is already a primitive recursive fast name in the completed
parameter circle.

One may make it a conventional real Cauchy name explicitly.
Choose the representative $t_0\in[0,1)$ and put $z_0=t_0$.
For each $n$, select the unique member
\[
 \Delta_n\in\{t_{n+1}-t_n-1,\ t_{n+1}-t_n,\ t_{n+1}-t_n+1\}
\]
of absolute value below $1/4$, and set $z_{n+1}=z_n+\Delta_n$.
The choice is a three-way rational comparison justified by
\eqref{eq:circlecauchy}. Moreover
$|\Delta_n|<2^{-(n+4)}$, so
\[
              |z_j-z_n|\le\sum_{i=n}^{j-1}2^{-(i+4)}
                         <2^{-(n+3)}\qquad(j>n).
\]
Thus $z_n$ is an explicitly primitive recursive Cauchy name;
its residue modulo one has name $t_n$. Let $t$ denote this
circle name. Forward continuity makes $\gamma(t_n)$ converge
to $\gamma(t)$ with an explicit rate, while
\eqref{eq:nearparameter} makes it converge to $x$.
For any required error choose
$n\ge\max(U(k+2),k+2)$ and combine the two bounds.
They show $\gamma(t)=x$ in the equality-of-names sense.

This constructs a preimage without compactness, an infinite
subsequence selection, or an unbounded search. Consequently
\[
 \Gamma=\{x:\text{every neighborhood of }x
                    \text{ contains rational points of }I\text{ and }O\}.
\]
Together with the local constancy and both-sides results, this
is the asserted constructive common-boundary statement.
All parts of Theorem~\ref{thm:main} are now proved.
